\documentclass[a4paper,12pt]{article}
\usepackage[utf8]{inputenc}
\usepackage[ngerman]{babel}
\usepackage{amsmath,amssymb}
\usepackage{geometry}
\geometry{margin=2.5cm}
\usepackage{enumitem}
\usepackage{booktabs}

\title{Übungsblatt 3 -- Lösungen \\ \large Linear Algebra}
\author{}
\date{}

\begin{document}
\maketitle

\noindent
\textbf{Erinnerung:} Zum Nachweis, dass eine Teilmenge $W$ ein Unterraum eines Vektorraums $V$ ist, prüfen wir drei Bedingungen:
\begin{enumerate}
    \item Der Nullvektor liegt in $W$: $\mathbf{0} \in W$.
    \item $W$ ist abgeschlossen unter Addition: $\mathbf{u}, \mathbf{v} \in W \Rightarrow \mathbf{u} + \mathbf{v} \in W$.
    \item $W$ ist abgeschlossen unter Skalarmultiplikation: $\mathbf{u} \in W,\; k \in \mathbb{R} \Rightarrow k\mathbf{u} \in W$.
\end{enumerate}

%% ============================================================
\subsection*{Aufgabe 10: Ist $W$ ein Unterraum von $V$?}
%% ============================================================

\begin{enumerate}[label=(\alph*)]

%% --- 10(a) ---
\item $V = \mathbb{R}^4$, \quad $W = \{(a,b,c,d) \in \mathbb{R}^4 : a + 2c = b - 3d\}$.

Die Bedingung lässt sich äquivalent als homogene lineare Gleichung schreiben:
\[
a - b + 2c + 3d = 0.
\]

\textbf{Nullvektor:}
Für $(0,0,0,0)$ gilt $0 - 0 + 2 \cdot 0 + 3 \cdot 0 = 0$. Also $\mathbf{0} \in W$. \checkmark

\textbf{Abgeschlossenheit unter Addition:}
Seien $\mathbf{u} = (a_1,b_1,c_1,d_1), \mathbf{v} = (a_2,b_2,c_2,d_2) \in W$.
Dann gilt $a_1 + 2c_1 = b_1 - 3d_1$ und $a_2 + 2c_2 = b_2 - 3d_2$.
Für $\mathbf{u} + \mathbf{v} = (a_1+a_2,\, b_1+b_2,\, c_1+c_2,\, d_1+d_2)$:
\[
(a_1+a_2) + 2(c_1+c_2) = (a_1+2c_1) + (a_2+2c_2) = (b_1-3d_1) + (b_2-3d_2) = (b_1+b_2) - 3(d_1+d_2).
\]
Also $\mathbf{u}+\mathbf{v} \in W$. \checkmark

\textbf{Abgeschlossenheit unter Skalarmultiplikation:}
Sei $\mathbf{u} = (a,b,c,d) \in W$ und $k \in \mathbb{R}$. Dann $a + 2c = b - 3d$ und:
\[
ka + 2(kc) = k(a+2c) = k(b-3d) = kb - 3(kd).
\]
Also $k\mathbf{u} \in W$. \checkmark

\textbf{$W$ ist ein Unterraum von $\mathbb{R}^4$.} $\square$

\bigskip

%% --- 10(b) ---
\item $V = \mathbb{R}^4$, \quad $W = \{(a,b,c,d) \in \mathbb{R}^4 : a + 2c = b - 3\}$.

\textbf{Nullvektor:}
Für $(0,0,0,0)$: $a + 2c = 0 + 0 = 0$ und $b - 3 = 0 - 3 = -3$.
Da $0 \neq -3$, gilt $\mathbf{0} \notin W$.

Da der Nullvektor nicht in $W$ liegt, ist die erste Bedingung verletzt.

\textbf{$W$ ist kein Unterraum von $\mathbb{R}^4$.} $\square$

\bigskip

%% --- 10(c) ---
\item $V = \mathbb{R}^3$, \quad $W = \{(a,b,c) \in \mathbb{R}^3 : a \leq c\}$.

\textbf{Nullvektor:} $(0,0,0)$: $0 \leq 0$ ist wahr, also $\mathbf{0} \in W$. \checkmark

\textbf{Abgeschlossenheit unter Skalarmultiplikation -- Gegenbeispiel:}
Wähle $\mathbf{u} = (1, 0, 2) \in W$, denn $1 \leq 2$. \checkmark

Für $k = -1$: $k\mathbf{u} = (-1)(1,0,2) = (-1, 0, -2)$.
Prüfe: $a = -1$, $c = -2$. Ist $-1 \leq -2$? \textbf{Nein}, denn $-1 > -2$.

Also ist $k\mathbf{u} \notin W$. Die Abgeschlossenheit unter Skalarmultiplikation ist verletzt.

\textbf{$W$ ist kein Unterraum von $\mathbb{R}^3$.} $\square$

\bigskip

%% --- 10(d) ---
\item $V = \mathbb{R}^3$, \quad $W = \{(a,b,c) \in \mathbb{R}^3 : |a| \leq |c|\}$.

\textbf{Nullvektor:} $(0,0,0)$: $|0| \leq |0|$ ist wahr, also $\mathbf{0} \in W$. \checkmark

\textbf{Abgeschlossenheit unter Addition -- Gegenbeispiel:}
Wähle $\mathbf{u} = (1, 0, -1) \in W$, denn $|1| = 1 \leq 1 = |-1|$. \checkmark

Wähle $\mathbf{v} = (1, 0, 1) \in W$, denn $|1| = 1 \leq 1 = |1|$. \checkmark

$\mathbf{u} + \mathbf{v} = (2, 0, 0)$.
Prüfe: $|2| = 2 \leq 0 = |0|$? \textbf{Nein}, denn $2 > 0$.

Also ist $\mathbf{u}+\mathbf{v} \notin W$. Die Abgeschlossenheit unter Addition ist verletzt.

\textbf{$W$ ist kein Unterraum von $\mathbb{R}^3$.} $\square$

\bigskip

%% --- 10(e) ---
\item $V = \mathbb{R}^3$, \quad $W = \{(a,b,c) \in \mathbb{R}^3 : ac \leq 0\}$.

\textbf{Nullvektor:} $(0,0,0)$: $0 \cdot 0 = 0 \leq 0$. Also $\mathbf{0} \in W$. \checkmark

\textbf{Abgeschlossenheit unter Addition -- Gegenbeispiel:}
Wähle $\mathbf{u} = (1, 0, 0) \in W$, denn $1 \cdot 0 = 0 \leq 0$. \checkmark

Wähle $\mathbf{v} = (0, 0, 1) \in W$, denn $0 \cdot 1 = 0 \leq 0$. \checkmark

$\mathbf{u} + \mathbf{v} = (1, 0, 1)$.
Prüfe: $1 \cdot 1 = 1 \leq 0$? \textbf{Nein}, denn $1 > 0$.

Also ist $\mathbf{u}+\mathbf{v} \notin W$. Die Abgeschlossenheit unter Addition ist verletzt.

\textbf{$W$ ist kein Unterraum von $\mathbb{R}^3$.} $\square$

\bigskip

%% --- 10(f) ---
\item $V = \mathbb{R}^3$, \quad $W = \{(a,b,c) \in \mathbb{R}^3 : a + b \neq c + 1\}$.

\textbf{Nullvektor:}
Für $(0,0,0)$: $a+b = 0$ und $c+1 = 1$. Da $0 \neq 1$, gilt $\mathbf{0} \in W$. \checkmark

\textbf{Abgeschlossenheit unter Addition -- Gegenbeispiel:}
Wähle $\mathbf{u} = (2, 0, 0) \in W$, denn $2 + 0 = 2$ und $0 + 1 = 1$, also $2 \neq 1$. \checkmark

Wähle $\mathbf{v} = (0, 0, 1) \in W$, denn $0 + 0 = 0$ und $1 + 1 = 2$, also $0 \neq 2$. \checkmark

$\mathbf{u} + \mathbf{v} = (2, 0, 1)$.
Prüfe: $a+b = 2+0 = 2$ und $c+1 = 1+1 = 2$. Es gilt $2 = 2$, also $a+b = c+1$.

Damit ist $\mathbf{u}+\mathbf{v} \notin W$. Die Abgeschlossenheit unter Addition ist verletzt.

\textbf{$W$ ist kein Unterraum von $\mathbb{R}^3$.} $\square$

\bigskip

%% --- 10(g) ---
\item $V = \mathbb{R}^4$, \quad $W = \{\ell \in \mathbb{R}^4 : \ell \text{ ist eine arithmetische Folge}\}$.

Eine arithmetische Folge in $\mathbb{R}^4$ hat die Form $(a,\, a+d,\, a+2d,\, a+3d)$ für $a, d \in \mathbb{R}$.
Schreiben wir $\ell = (\ell_1, \ell_2, \ell_3, \ell_4)$, so ist $\ell$ genau dann arithmetisch, wenn die aufeinanderfolgenden Differenzen gleich sind:
\[
\ell_2 - \ell_1 = \ell_3 - \ell_2 = \ell_4 - \ell_3.
\]
Dies ist äquivalent zu den beiden homogenen linearen Gleichungen:
\[
\ell_1 - 2\ell_2 + \ell_3 = 0 \quad \text{und} \quad \ell_2 - 2\ell_3 + \ell_4 = 0.
\]

\textbf{Nullvektor:}
$(0,0,0,0)$: arithmetisch mit $a=0, d=0$. Beide Gleichungen: $0-0+0=0$. Also $\mathbf{0} \in W$. \checkmark

\textbf{Abgeschlossenheit unter Addition:}
Seien $\mathbf{u} = (a_1, a_1+d_1, a_1+2d_1, a_1+3d_1)$ und $\mathbf{v} = (a_2, a_2+d_2, a_2+2d_2, a_2+3d_2) \in W$.
\begin{align*}
\mathbf{u}+\mathbf{v} &= \big(a_1+a_2,\; (a_1+a_2)+(d_1+d_2),\; (a_1+a_2)+2(d_1+d_2),\; (a_1+a_2)+3(d_1+d_2)\big).
\end{align*}
Dies ist eine arithmetische Folge mit Anfangswert $a_1+a_2$ und Differenz $d_1+d_2$.
Also $\mathbf{u}+\mathbf{v} \in W$. \checkmark

\textbf{Abgeschlossenheit unter Skalarmultiplikation:}
Sei $\mathbf{u} = (a, a+d, a+2d, a+3d) \in W$ und $k \in \mathbb{R}$.
\[
k\mathbf{u} = (ka,\; ka+kd,\; ka+2kd,\; ka+3kd).
\]
Dies ist eine arithmetische Folge mit Anfangswert $ka$ und Differenz $kd$.
Also $k\mathbf{u} \in W$. \checkmark

\textbf{$W$ ist ein Unterraum von $\mathbb{R}^4$.} $\square$

\bigskip

%% --- 10(h) ---
\item $V = \mathbb{R}^4$, \quad $W = \{\ell \in \mathbb{R}^4 : \ell \text{ ist eine geometrische Folge}\}$.

Eine geometrische Folge hat die Form $(b, bq, bq^2, bq^3)$ für $b, q \in \mathbb{R}$.

\textbf{Abgeschlossenheit unter Addition -- Gegenbeispiel:}
Wähle $\mathbf{u} = (1, 2, 4, 8) \in W$ (geometrisch mit $b=1, q=2$). \checkmark

Wähle $\mathbf{v} = (1, 3, 9, 27) \in W$ (geometrisch mit $b=1, q=3$). \checkmark

$\mathbf{u}+\mathbf{v} = (2, 5, 13, 35)$.

Falls dies eine geometrische Folge wäre, müsste der Quotient konstant sein:
\[
\frac{5}{2} = 2{,}5 \quad \text{aber} \quad \frac{13}{5} = 2{,}6.
\]
Da $2{,}5 \neq 2{,}6$, ist $(2, 5, 13, 35)$ keine geometrische Folge.

Also ist $\mathbf{u}+\mathbf{v} \notin W$. Die Abgeschlossenheit unter Addition ist verletzt.

\textbf{$W$ ist kein Unterraum von $\mathbb{R}^4$.} $\square$

\bigskip

%% --- 10(i) ---
\item $V = \mathbb{R}^3$, \quad $W = \mathbb{R}^2$.

$\mathbb{R}^2$ besteht aus geordneten Paaren $(a,b)$ und $\mathbb{R}^3$ aus geordneten Tripeln $(a,b,c)$. Dies sind verschiedene mathematische Objekte: Ein Paar $(a,b)$ ist kein Element von $\mathbb{R}^3$, da die Elemente von $\mathbb{R}^3$ stets drei Komponenten haben.

Da $\mathbb{R}^2 \not\subseteq \mathbb{R}^3$, kann $W = \mathbb{R}^2$ kein Unterraum von $V = \mathbb{R}^3$ sein. Die Frage ist bereits auf der Ebene der Mengenlehre nicht sinnvoll: $W$ ist keine Teilmenge von $V$.

\textbf{$W$ ist kein Unterraum von $\mathbb{R}^3$.} $\square$

\end{enumerate}

\bigskip
\hrule
\bigskip

%% ============================================================
\subsection*{Aufgabe 11: Ist $W$ ein Unterraum von $\mathbb{R}^{n \times n}$ $(n > 1)$?}
%% ============================================================

\begin{enumerate}[label=(\alph*)]

%% --- 11(a) ---
\item $W = \{A \in \mathbb{R}^{n \times n} : A^\top = -A\}$ \quad (schiefsymmetrische Matrizen).

\textbf{Nullvektor:}
Die Nullmatrix $0$ erfüllt $0^\top = 0 = -0$. Also $0 \in W$. \checkmark

\textbf{Abgeschlossenheit unter Addition:}
Seien $A, B \in W$, also $A^\top = -A$ und $B^\top = -B$. Dann:
\[
(A+B)^\top = A^\top + B^\top = (-A) + (-B) = -(A+B).
\]
Also $A+B \in W$. \checkmark

\textbf{Abgeschlossenheit unter Skalarmultiplikation:}
Sei $A \in W$ und $k \in \mathbb{R}$. Dann $A^\top = -A$ und:
\[
(kA)^\top = kA^\top = k(-A) = -(kA).
\]
Also $kA \in W$. \checkmark

\textbf{$W$ ist ein Unterraum von $\mathbb{R}^{n \times n}$.} $\square$

\bigskip

%% --- 11(b) ---
\item $W = \{A \in \mathbb{R}^{n \times n} : A^2 = 0\}$.

\textbf{Abgeschlossenheit unter Addition -- Gegenbeispiel:}
Wähle $n = 2$ und:
\[
A = \begin{pmatrix} 0 & 1 \\ 0 & 0 \end{pmatrix}, \qquad
B = \begin{pmatrix} 0 & 0 \\ 1 & 0 \end{pmatrix}.
\]

Prüfe $A^2$:
\[
A^2 = \begin{pmatrix} 0 & 1 \\ 0 & 0 \end{pmatrix} \begin{pmatrix} 0 & 1 \\ 0 & 0 \end{pmatrix} = \begin{pmatrix} 0 & 0 \\ 0 & 0 \end{pmatrix} = 0. \checkmark
\]

Prüfe $B^2$:
\[
B^2 = \begin{pmatrix} 0 & 0 \\ 1 & 0 \end{pmatrix} \begin{pmatrix} 0 & 0 \\ 1 & 0 \end{pmatrix} = \begin{pmatrix} 0 & 0 \\ 0 & 0 \end{pmatrix} = 0. \checkmark
\]

Aber:
\[
A + B = \begin{pmatrix} 0 & 1 \\ 1 & 0 \end{pmatrix}, \qquad
(A+B)^2 = \begin{pmatrix} 0 & 1 \\ 1 & 0 \end{pmatrix}\begin{pmatrix} 0 & 1 \\ 1 & 0 \end{pmatrix} = \begin{pmatrix} 1 & 0 \\ 0 & 1 \end{pmatrix} = I \neq 0.
\]

Also ist $A+B \notin W$. Die Abgeschlossenheit unter Addition ist verletzt.

\textbf{$W$ ist kein Unterraum von $\mathbb{R}^{n \times n}$.} $\square$

\bigskip

%% --- 11(c) ---
\item $W = \{A \in \mathbb{R}^{n \times n} : A \cdot \mathbf{v} = \mathbf{0}\}$, wobei $\mathbf{v} = (1, 2, \ldots, n)^\top$.

Dies ist der Kern der linearen Abbildung $\varphi \colon \mathbb{R}^{n \times n} \to \mathbb{R}^n$, $\varphi(A) = A\mathbf{v}$.

\textbf{Nullvektor:}
$0 \cdot \mathbf{v} = \mathbf{0}$. Also $0 \in W$. \checkmark

\textbf{Abgeschlossenheit unter Addition:}
Seien $A, B \in W$, also $A\mathbf{v} = \mathbf{0}$ und $B\mathbf{v} = \mathbf{0}$. Dann:
\[
(A+B)\mathbf{v} = A\mathbf{v} + B\mathbf{v} = \mathbf{0} + \mathbf{0} = \mathbf{0}.
\]
Also $A+B \in W$. \checkmark

\textbf{Abgeschlossenheit unter Skalarmultiplikation:}
Sei $A \in W$ und $k \in \mathbb{R}$. Dann:
\[
(kA)\mathbf{v} = k(A\mathbf{v}) = k \cdot \mathbf{0} = \mathbf{0}.
\]
Also $kA \in W$. \checkmark

\textbf{$W$ ist ein Unterraum von $\mathbb{R}^{n \times n}$.} $\square$

\bigskip

%% --- 11(d) ---
\item $W = \{A \in \mathbb{R}^{n \times n} : A \text{ ist in Zeilenstufenform}\}$.

\textbf{Abgeschlossenheit unter Addition -- Gegenbeispiel:}
Wähle $n = 2$ und:
\[
A = \begin{pmatrix} 1 & 0 \\ 0 & 1 \end{pmatrix}, \qquad
B = \begin{pmatrix} -1 & 0 \\ 0 & 0 \end{pmatrix}.
\]

$A$ ist in Zeilenstufenform (Pivots in Position $(1,1)$ und $(2,2)$). \checkmark

$B$ ist in Zeilenstufenform (Pivot in Position $(1,1)$, zweite Zeile ist Nullzeile). \checkmark

Aber:
\[
A + B = \begin{pmatrix} 0 & 0 \\ 0 & 1 \end{pmatrix}.
\]

Die erste Zeile ist eine Nullzeile, aber die zweite Zeile ist keine Nullzeile. In Zeilenstufenform müssen alle Nullzeilen unterhalb der Nicht-Nullzeilen stehen. Diese Bedingung ist verletzt.

Also ist $A+B \notin W$. Die Abgeschlossenheit unter Addition ist verletzt.

\textbf{$W$ ist kein Unterraum von $\mathbb{R}^{n \times n}$.} $\square$

\bigskip

%% --- 11(e) ---
\item $W = \{A \in \mathbb{R}^{n \times n} : KAM = A^\top\}$, wobei $K, M \in \mathbb{R}^{n \times n}$ fest gewählt.

\textbf{Nullvektor:}
$K \cdot 0 \cdot M = 0 = 0^\top$. Also $0 \in W$. \checkmark

\textbf{Abgeschlossenheit unter Addition:}
Seien $A, B \in W$, also $KAM = A^\top$ und $KBM = B^\top$. Dann:
\[
K(A+B)M = KAM + KBM = A^\top + B^\top = (A+B)^\top.
\]
Also $A+B \in W$. \checkmark

\textbf{Abgeschlossenheit unter Skalarmultiplikation:}
Sei $A \in W$ und $k \in \mathbb{R}$. Dann:
\[
K(kA)M = k(KAM) = kA^\top = (kA)^\top.
\]
Also $kA \in W$. \checkmark

\textbf{$W$ ist ein Unterraum von $\mathbb{R}^{n \times n}$.} $\square$

\bigskip

%% --- 11(f) ---
\item $W = \{A \in \mathbb{R}^{n \times n} : \operatorname{Spur}(A) = 0\}$.

Erinnerung: $\operatorname{Spur}(A) = \sum_{i=1}^{n} a_{ii}$.

\textbf{Nullvektor:}
$\operatorname{Spur}(0) = 0$. Also $0 \in W$. \checkmark

\textbf{Abgeschlossenheit unter Addition:}
Seien $A, B \in W$, also $\operatorname{Spur}(A) = 0$ und $\operatorname{Spur}(B) = 0$. Dann:
\[
\operatorname{Spur}(A+B) = \operatorname{Spur}(A) + \operatorname{Spur}(B) = 0 + 0 = 0.
\]
Also $A+B \in W$. \checkmark

\textbf{Abgeschlossenheit unter Skalarmultiplikation:}
Sei $A \in W$ und $k \in \mathbb{R}$. Dann:
\[
\operatorname{Spur}(kA) = k \cdot \operatorname{Spur}(A) = k \cdot 0 = 0.
\]
Also $kA \in W$. \checkmark

\textbf{$W$ ist ein Unterraum von $\mathbb{R}^{n \times n}$.} $\square$

\end{enumerate}

\bigskip
\hrule
\bigskip

%% ============================================================
\subsection*{Aufgabe 12: Ist $W$ ein Unterraum von $\mathbb{R}[x]$?}
%% ============================================================

Hier ist $\mathbb{R}[x]$ der Vektorraum aller Polynome mit reellen Koeffizienten.

\begin{enumerate}[label=(\alph*)]

%% --- 12(a) ---
\item $W = \{p(x) \in \mathbb{R}[x] : p(1) = 0\}$.

\textbf{Nullvektor:}
Das Nullpolynom $p(x) = 0$ erfüllt $p(1) = 0$. Also $0 \in W$. \checkmark

\textbf{Abgeschlossenheit unter Addition:}
Seien $p, q \in W$, also $p(1) = 0$ und $q(1) = 0$. Dann:
\[
(p+q)(1) = p(1) + q(1) = 0 + 0 = 0.
\]
Also $p+q \in W$. \checkmark

\textbf{Abgeschlossenheit unter Skalarmultiplikation:}
Sei $p \in W$ und $k \in \mathbb{R}$. Dann:
\[
(kp)(1) = k \cdot p(1) = k \cdot 0 = 0.
\]
Also $kp \in W$. \checkmark

\textbf{$W$ ist ein Unterraum von $\mathbb{R}[x]$.} $\square$

\bigskip

%% --- 12(b) ---
\item $W = \{p(x) \in \mathbb{R}[x] : p(1) = 0 \text{ und } p(2) = 0\}$.

\textbf{Nullvektor:}
Das Nullpolynom erfüllt $0(1) = 0$ und $0(2) = 0$. Also $0 \in W$. \checkmark

\textbf{Abgeschlossenheit unter Addition:}
Seien $p, q \in W$. Dann $p(1) = q(1) = 0$ und $p(2) = q(2) = 0$.
\[
(p+q)(1) = p(1) + q(1) = 0 + 0 = 0, \qquad (p+q)(2) = p(2) + q(2) = 0 + 0 = 0.
\]
Also $p+q \in W$. \checkmark

\textbf{Abgeschlossenheit unter Skalarmultiplikation:}
Sei $p \in W$ und $k \in \mathbb{R}$.
\[
(kp)(1) = k \cdot p(1) = 0, \qquad (kp)(2) = k \cdot p(2) = 0.
\]
Also $kp \in W$. \checkmark

\textbf{$W$ ist ein Unterraum von $\mathbb{R}[x]$.} $\square$

\bigskip

%% --- 12(c) ---
\item $W = \{p(x) \in \mathbb{R}[x] : p(1) \cdot p(2) = 0\}$.

\textbf{Nullvektor:} $0(1) \cdot 0(2) = 0 \cdot 0 = 0$. Also $0 \in W$. \checkmark

\textbf{Abgeschlossenheit unter Addition -- Gegenbeispiel:}
Wähle $p(x) = x - 1$ und $q(x) = x - 2$.

$p(1) = 0$, $p(2) = 1$: $p(1) \cdot p(2) = 0 \cdot 1 = 0$. Also $p \in W$. \checkmark

$q(1) = -1$, $q(2) = 0$: $q(1) \cdot q(2) = (-1) \cdot 0 = 0$. Also $q \in W$. \checkmark

$(p+q)(x) = 2x - 3$. Auswertung: $(p+q)(1) = -1$, $(p+q)(2) = 1$.
\[
(p+q)(1) \cdot (p+q)(2) = (-1) \cdot 1 = -1 \neq 0.
\]

Also ist $p+q \notin W$. Die Abgeschlossenheit unter Addition ist verletzt.

\textbf{$W$ ist kein Unterraum von $\mathbb{R}[x]$.} $\square$

\bigskip

%% --- 12(d) ---
\item $W = \{p(x) \in \mathbb{R}[x] : p(1) \leq p(2)\}$.

\textbf{Nullvektor:} $0(1) = 0 \leq 0 = 0(2)$. Also $0 \in W$. \checkmark

\textbf{Abgeschlossenheit unter Skalarmultiplikation -- Gegenbeispiel:}
Wähle $p(x) = x$. Dann $p(1) = 1 \leq 2 = p(2)$, also $p \in W$. \checkmark

Für $k = -1$: $(-p)(x) = -x$. Auswertung: $(-p)(1) = -1$, $(-p)(2) = -2$.
Ist $-1 \leq -2$? \textbf{Nein}, denn $-1 > -2$.

Also ist $kp \notin W$. Die Abgeschlossenheit unter Skalarmultiplikation ist verletzt.

\textbf{$W$ ist kein Unterraum von $\mathbb{R}[x]$.} $\square$

\bigskip

%% --- 12(e) ---
\item $W = \{p(x) \in \mathbb{R}[x] : p(1) = p(2) + p(3)\}$.

Die Bedingung $p(1) - p(2) - p(3) = 0$ ist eine lineare Bedingung an $p$ (die Auswertung an einer Stelle ist eine lineare Abbildung).

\textbf{Nullvektor:} $0(1) = 0 = 0 + 0 = 0(2) + 0(3)$. Also $0 \in W$. \checkmark

\textbf{Abgeschlossenheit unter Addition:}
Seien $p, q \in W$, also $p(1) = p(2) + p(3)$ und $q(1) = q(2) + q(3)$. Dann:
\begin{align*}
(p+q)(1) &= p(1) + q(1) = \big(p(2) + p(3)\big) + \big(q(2) + q(3)\big) \\
         &= \big(p(2) + q(2)\big) + \big(p(3) + q(3)\big) = (p+q)(2) + (p+q)(3).
\end{align*}
Also $p+q \in W$. \checkmark

\textbf{Abgeschlossenheit unter Skalarmultiplikation:}
Sei $p \in W$ und $k \in \mathbb{R}$. Dann:
\[
(kp)(1) = k \cdot p(1) = k\big(p(2) + p(3)\big) = kp(2) + kp(3) = (kp)(2) + (kp)(3).
\]
Also $kp \in W$. \checkmark

\textbf{$W$ ist ein Unterraum von $\mathbb{R}[x]$.} $\square$

\bigskip

%% --- 12(f) ---
\item $W = \{p(x) \in \mathbb{R}[x] : p'(1) = p(0)\}$.

Die Bedingung $p'(1) - p(0) = 0$ ist linear, da sowohl die Auswertung $p \mapsto p(0)$ als auch $p \mapsto p'(1)$ lineare Abbildungen sind.

\textbf{Nullvektor:}
$0'(x) = 0$, also $0'(1) = 0 = 0(0)$. Also $0 \in W$. \checkmark

\textbf{Abgeschlossenheit unter Addition:}
Seien $p, q \in W$, also $p'(1) = p(0)$ und $q'(1) = q(0)$. Dann:
\[
(p+q)'(1) = p'(1) + q'(1) = p(0) + q(0) = (p+q)(0).
\]
Also $p+q \in W$. \checkmark

\textbf{Abgeschlossenheit unter Skalarmultiplikation:}
Sei $p \in W$ und $k \in \mathbb{R}$. Dann:
\[
(kp)'(1) = k \cdot p'(1) = k \cdot p(0) = (kp)(0).
\]
Also $kp \in W$. \checkmark

\textbf{$W$ ist ein Unterraum von $\mathbb{R}[x]$.} $\square$

\end{enumerate}

\bigskip
\hrule
\bigskip

%% ============================================================
\subsection*{Aufgabe 13: Ist $W$ ein Unterraum von $\operatorname{FUNC}(\mathbb{R})$?}
%% ============================================================

Hier ist $\operatorname{FUNC}(\mathbb{R})$ der Vektorraum aller Funktionen $f \colon \mathbb{R} \to \mathbb{R}$.

\begin{enumerate}[label=(\alph*)]

%% --- 13(a) ---
\item $W = \{f \in \operatorname{FUNC}(\mathbb{R}) : f(-x) = -f(x) \text{ für alle } x \in \mathbb{R}\}$ \quad (ungerade Funktionen).

\textbf{Nullvektor:}
Die Nullfunktion $f(x) = 0$ erfüllt $f(-x) = 0 = -0 = -f(x)$ für alle $x$. Also $0 \in W$. \checkmark

\textbf{Abgeschlossenheit unter Addition:}
Seien $f, g \in W$. Dann für alle $x \in \mathbb{R}$:
\[
(f+g)(-x) = f(-x) + g(-x) = -f(x) + (-g(x)) = -(f(x) + g(x)) = -(f+g)(x).
\]
Also $f+g \in W$. \checkmark

\textbf{Abgeschlossenheit unter Skalarmultiplikation:}
Sei $f \in W$ und $k \in \mathbb{R}$. Dann für alle $x \in \mathbb{R}$:
\[
(kf)(-x) = k \cdot f(-x) = k \cdot (-f(x)) = -(kf)(x).
\]
Also $kf \in W$. \checkmark

\textbf{$W$ ist ein Unterraum von $\operatorname{FUNC}(\mathbb{R})$.} $\square$

\bigskip

%% --- 13(b) ---
\item $W = \{f \in \operatorname{FUNC}(\mathbb{R}) : f(x) \neq 0 \text{ für alle } x \neq 0\}$.

\textbf{Nullvektor:}
Die Nullfunktion $f(x) = 0$ erfüllt $f(1) = 0$, obwohl $1 \neq 0$. Damit ist $f \notin W$.

Da der Nullvektor nicht in $W$ liegt, ist die erste Bedingung verletzt.

\textit{Alternativ} (Abgeschlossenheit unter Addition): Die Funktionen $f(x) = x$ und $g(x) = -x$ liegen beide in $W$ (für $x \neq 0$ gilt $f(x) = x \neq 0$ und $g(x) = -x \neq 0$). Aber $(f+g)(x) = 0$ für alle $x$, insbesondere $(f+g)(1) = 0$, also $f+g \notin W$.

\textbf{$W$ ist kein Unterraum von $\operatorname{FUNC}(\mathbb{R})$.} $\square$

\bigskip

%% --- 13(c) ---
\item $W = \{f \in \operatorname{FUNC}(\mathbb{R}) : f(x) = f(x+1) \text{ für alle } x \in \mathbb{R}\}$ \quad (1-periodische Funktionen).

\textbf{Nullvektor:}
$0(x) = 0 = 0(x+1)$ für alle $x$. Also $0 \in W$. \checkmark

\textbf{Abgeschlossenheit unter Addition:}
Seien $f, g \in W$, also $f(x) = f(x+1)$ und $g(x) = g(x+1)$ für alle $x$. Dann:
\[
(f+g)(x) = f(x) + g(x) = f(x+1) + g(x+1) = (f+g)(x+1).
\]
Also $f+g \in W$. \checkmark

\textbf{Abgeschlossenheit unter Skalarmultiplikation:}
Sei $f \in W$ und $k \in \mathbb{R}$. Dann für alle $x$:
\[
(kf)(x) = k \cdot f(x) = k \cdot f(x+1) = (kf)(x+1).
\]
Also $kf \in W$. \checkmark

\textbf{$W$ ist ein Unterraum von $\operatorname{FUNC}(\mathbb{R})$.} $\square$

\bigskip

%% --- 13(d) ---
\item $W = \{f \in \operatorname{FUNC}(\mathbb{R}) : f(x) + 1 = f(x+1) \text{ für alle } x \in \mathbb{R}\}$.

\textbf{Nullvektor:}
Die Nullfunktion $f(x) = 0$ erfüllt $f(x) + 1 = 0 + 1 = 1$ und $f(x+1) = 0$.
Da $1 \neq 0$, gilt $f \notin W$.

Da der Nullvektor nicht in $W$ liegt, ist die erste Bedingung verletzt.

\textit{Bemerkung:} Ein Beispiel einer Funktion in $W$ wäre $f(x) = x$, denn $x + 1 = (x+1)$. Aber das Nullelement des Vektorraums liegt nicht in $W$.

\textbf{$W$ ist kein Unterraum von $\operatorname{FUNC}(\mathbb{R})$.} $\square$

\end{enumerate}

\end{document}
